% ============================================================
% 双栏论文版式（期刊风 / 参考件那种排法）
%
% 必须同时传 classoption，否则还是单栏：
%   pandoc src/report.md -o out.pdf --pdf-engine=xelatex \
%     --include-in-header=templates/header-paper.tex \
%     -V documentclass=article -V classoption=twocolumn \
%     -V fontsize=10pt -V papersize=a4
%
% 与 header-cn.tex 的差别：双栏、10pt、思源宋体、栏间距 6mm。
% 双栏下每行只有 22–24 个汉字，再大字号就挤；别用 11pt。
%
% 跨栏的表和图用 table* / figure*（只能落页顶，不能 [H]）。
% ============================================================

% ---- 字体：中文思源黑/宋，西文 Times ----
% 用 Noto Serif SC，不要用 Source Han Serif SC：后者在本机安装的 Heavy 字重
% 缺 、。 等 CJK 标点，一旦出现加粗中文就静默吞标点（实测 15 个缺字）。
\usepackage{xeCJK}
\setCJKmainfont[AutoFakeBold=2]{Noto Serif SC}
\setCJKsansfont[AutoFakeBold=2]{SimHei}
\setCJKmonofont{SimSun}
\setmainfont{Times New Roman}
\setsansfont{Arial}
\setmonofont{Consolas}

% ---- 版面：双栏 ----
\usepackage[a4paper,top=2cm,bottom=2cm,left=2cm,right=2cm,columnsep=6mm]{geometry}
\linespread{1.15}
\setlength{\parindent}{2em}
\setlength{\parskip}{0pt}
\setlength{\emergencystretch}{2em}
\widowpenalty=10000
\clubpenalty=10000

% ---- 标题：双栏下字号收紧一档 ----
\usepackage{titlesec}
\titleformat{\section}{\sffamily\fontsize{12pt}{15pt}\selectfont\bfseries}{\thesection}{0.6em}{}
\titlespacing*{\section}{0pt}{10pt plus 2pt}{5pt plus 1pt}
\titleformat{\subsection}{\sffamily\fontsize{10pt}{13pt}\selectfont\bfseries}{\thesubsection}{0.5em}{}
\titlespacing*{\subsection}{0pt}{8pt plus 1pt}{3pt}
\titleformat{\subsubsection}{\sffamily\fontsize{10pt}{13pt}\selectfont\bfseries}{\thesubsubsection}{0.5em}{}
\titlespacing*{\subsubsection}{0pt}{6pt plus 1pt}{2pt}

% ---- 交叉引用宏 ----
\newcommand{\tabref}[1]{表~\ref{#1}}
\newcommand{\figref}[1]{图~\ref{#1}}

% ---- 题注：双栏下再小一号 ----
\usepackage[font=footnotesize,labelsep=quad,skip=4pt,justification=centering]{caption}
\captionsetup[table]{position=top}
\renewcommand{\tablename}{表}
\renewcommand{\figurename}{图}

% ---- 三线表 ----
\usepackage{booktabs}
\usepackage{tabularx}
\usepackage{array}
\setlength{\heavyrulewidth}{1.5pt}
\setlength{\lightrulewidth}{0.75pt}
\setlength{\cmidrulewidth}{0.5pt}
\renewcommand{\arraystretch}{1.05}
\setlength{\tabcolsep}{4pt}
\newcolumntype{L}[1]{>{\raggedright\arraybackslash}p{#1}}

% ---- 图 ----
\usepackage{graphicx}
\usepackage{float}
\setlength{\textfloatsep}{10pt plus 2pt minus 2pt}
\setlength{\intextsep}{6pt plus 1pt}
\renewcommand{\topfraction}{0.92}
\renewcommand{\bottomfraction}{0.8}
\renewcommand{\textfraction}{0.08}
\renewcommand{\floatpagefraction}{0.75}
% 纯浮动页的胶水默认可伸缩，会把这些页的浮动体撑到页面中间，看起来"松散"。
% 顶部不伸缩、余量全堆到底部，浮动体就从页顶往下排。
\makeatletter
\setlength{\@fptop}{0pt}
\setlength{\@fpsep}{12pt}
\setlength{\@fpbot}{0pt plus 1fil}
\makeatother

% ---- 页眉页脚 ----
\usepackage{fancyhdr}
\providecommand{\doctitle}{}
\providecommand{\docauthor}{}
\pagestyle{fancy}
\fancyhf{}
\fancyhead[L]{\footnotesize\sffamily\doctitle}
\fancyhead[R]{\footnotesize\sffamily\docauthor}
\fancyfoot[C]{\footnotesize\thepage}
\renewcommand{\headrulewidth}{0.4pt}
\fancypagestyle{plain}{\fancyhf{}\fancyfoot[C]{\footnotesize\thepage}%
  \renewcommand{\headrulewidth}{0pt}}
\setlength{\headheight}{15pt}

% ---- 链接 ----
\usepackage[unicode,hidelinks]{hyperref}
\AtBeginDocument{\hypersetup{pdftitle={\doctitle},pdfauthor={\docauthor}}}
